\documentclass[11pt,a4paper]{article}
\usepackage[margin=2.15cm]{geometry}
\usepackage[spanish,es-minimal,es-nodecimaldot]{babel}
\usepackage{graphicx,booktabs,xcolor,amsmath,hyperref}
\usepackage[font=small,labelfont=bf,hypcap=false]{caption}

\definecolor{eva}{HTML}{67419A}
\definecolor{tinta}{HTML}{203438}
\hypersetup{colorlinks,linkcolor=eva,urlcolor=eva,pdftitle={Rendimiento agrícola en el Cesar}}
\graphicspath{{../figuras/}{documentos/figuras/}}
\setlength{\parindent}{0pt}
\setlength{\parskip}{5pt}
\renewcommand{\baselinestretch}{1.06}
\newcommand{\autores}{Sebastián Mindiola · William Andres Velasquez Ruiz · Jhon Eduardo Meneses Viecco}
\newcommand{\figura}[3]{%
  \begin{center}
    \begin{minipage}{\linewidth}
      \centering\includegraphics[width=#1]{#2}\captionof{figure}{#3}
    \end{minipage}
  \end{center}}

\begin{document}
\color{tinta}
{\color{eva}\small\textbf{UNIVERSIDAD NACIONAL DE COLOMBIA}\par}
\vspace{.25cm}
{\huge\bfseries Rendimiento agrícola en el Cesar\par}
{\large Evaluaciones Agropecuarias Municipales, 2007--2018\par}
\vspace{.3cm}
{\small \autores\par
Programación en Lenguajes Estadísticos (2026-2) · Docente Yoselin Nieto G.\par
25 de noviembre de 2026\par}

\section*{Resumen ejecutivo}
Estudiamos cómo varía el rendimiento agrícola entre cultivos, municipios y años.
El flujo en R y Python identificó 2.134 registros anuales para los diez cultivos principales,
que representan el 92,8\,\% del área sembrada. Marcamos seis rendimientos fuera de las cercas
exteriores del rango intercuartílico y repetimos el análisis con esas observaciones.

Las diferencias entre las medias de los cultivos reúnen el 90,59\,\% de la variabilidad
del rendimiento en t/ha. Dentro de los cultivos, la separación por municipio reúne
la mayor fracción entre las variables examinadas. La cobertura varía entre años y cultivos,
por lo que las trayectorias departamentales reflejan los municipios que reportaron cada año.

\section{Primera inspección}
El \href{https://www.datos.gov.co/d/2pnw-mmge}{registro EVA}
contiene 206.068 filas nacionales. El subconjunto del Cesar reúne 4.887 registros,
25 municipios y 50 cultivos. Cada fila identifica cultivo, municipio, año,
periodo y sistema productivo, con hectáreas sembradas y cosechadas, producción en toneladas y
rendimiento reportado en t/ha. El periodo permite distinguir los registros anuales de los
semestrales. En particular, \texttt{2016A} y \texttt{2016B} corresponden a los dos semestres de 2016.

Una vez convertidas las cantidades del CSV, \texttt{head()} devuelve las primeras filas del
subconjunto del Cesar en el orden del archivo. El cuadro recoge las columnas que usamos
en el análisis. Estas cinco filas corresponden a aguacate y el periodo \texttt{2007}
identifica un registro anual.

\begin{center}
\begin{minipage}{\linewidth}
\centering\small
\begin{tabular}{lrrrrr}
\toprule
Municipio & Periodo & Sembrada (ha) & Cosechada (ha) & Producción (t) & Rend. (t/ha)\\
\midrule
Valledupar          & 2007 & 713 & 713 & 5.704 & 8,00\\
La Jagua de Ibirico & 2007 & 430 & 430 & 2.150 & 5,00\\
La Paz              & 2007 & 300 & 300 & 3.000 & 10,00\\
Pueblo Bello        & 2007 & 105 & 103 & 830 & 8,06\\
San Martín          & 2007 & 90 & 80 & 640 & 8,00\\
\bottomrule
\end{tabular}
\captionof{table}{Primeras cinco filas del Cesar tras la importación. El rendimiento de la última columna es el que reporta EVA.}
\end{minipage}
\end{center}
En San Martín se sembraron 90 ha de aguacate y se cosecharon 80. Las 640 t producidas
equivalen a \(640/80=8\) t por hectárea cosechada. Al extender la revisión a los 4.887 registros
del Cesar, el cuadro muestra los límites de cada variable numérica y el número de filas sin valor.

\begin{center}
\begin{minipage}{\linewidth}
\centering\small
\begin{tabular}{lrrr}
\toprule
Variable & Mínimo & Máximo & Faltantes\\
\midrule
Área sembrada (ha) & 1 & 18.000 & 0\\
Área cosechada (ha) & 0 & 16.000 & 0\\
Producción (t) & 0 & 51.000 & 0\\
Rendimiento reportado (t/ha) & 0,20 & 184 & 53\\
\bottomrule
\end{tabular}
\captionof{table}{Rangos y valores faltantes en el subconjunto del Cesar antes de limpiar y agrupar. Los extremos del rango se calculan sobre los valores presentes.}
\end{minipage}
\end{center}

Los 53 rendimientos faltantes representan el 1,08\,\% de las filas. En todos esos casos
EVA reportó cero hectáreas cosechadas y cero toneladas de producción. Las otras 16 columnas
tienen datos en todas las filas.

Los mínimos y máximos indican los límites de la base. Los histogramas muestran cuántos
registros se concentran cerca de un extremo o se reparten a lo largo del rango. Cada panel
corresponde a una de las cuatro variables del cuadro y la altura de cada barra cuenta los
registros que caen en ese intervalo.

El área sembrada tiene una mediana de 138 ha y alcanza 18.000 ha. En un eje lineal, unos
pocos registros grandes comprimen la mayor parte de las barras cerca de cero. Usamos
el logaritmo decimal para acortar esa cola. Sumamos 1 dentro del logaritmo para conservar
las filas con valor cero, pues \(\log_{10}(1+0)=0\). Para cada valor original \(x_i\), la posición es
\[
u_i=\log_{10}(1+x_i).
\]
Por ejemplo, las 713 ha sembradas de Valledupar caen cerca de
\(\log_{10}(714)=2{,}85\) en el primer panel. La transformación fija la posición de cada
observación en el gráfico. Las cantidades originales se usan al sumar áreas y producción.

\figura{.91\linewidth}{eva_Py_01_inspeccion.png}{Distribución de las áreas sembrada y cosechada, la producción
y el rendimiento reportado en Cesar. El eje horizontal usa la escala transformada y el vertical cuenta registros.}

\section{Limpieza y construcción de la tabla anual}

Leímos el CSV como texto y convertimos las cantidades con punto de miles y coma decimal.
Estandarizamos los textos y conservamos los códigos municipales como identificadores.
La revisión encontró cero duplicados y cero filas con cantidades incompatibles en las áreas
y la producción. Hubo 59 discrepancias superiores a la tolerancia
\(0{,}05\lvert y_{\text{reportado}}\rvert+0{,}01\) entre el rendimiento reportado y el cociente recalculado.
Las marcamos para consulta y usamos las cantidades aditivas para calcular el rendimiento anual.

Conservamos las 9.580 hectáreas sembradas de los registros sin cosecha para medir la superficie
dedicada a la siembra. Al agregar los datos, sumamos las hectáreas de los cultivos y años elegibles.
El año 2006 solo presenta registros del periodo B en el archivo de Cesar y queda fuera
de la agregación anual. Para 2007--2018 sumamos las cantidades por municipio, cultivo y año,
y calculamos
\[
Y_{mct}=\frac{\sum_{i\in(m,c,t)}\mathrm{producci\acute on}_i}
                  {\sum_{i\in(m,c,t)}\mathrm{\acute area\ cosechada}_i}.
\]
Cuando el denominador es cero, el rendimiento queda indefinido. El criterio de selección de
los diez cultivos principales fue el área sembrada acumulada en los años completos.
De las 2.134 filas anuales, 11 carecen de cosecha y 2.123 tienen rendimiento definido.


\subsection*{De dos registros de arroz a una observación anual}
En Curumaní, EVA reportó 131 t de arroz en 88 ha durante 2016A. En 2016B
reportó 25 t en 20 ha. La fila anual suma por separado producción y superficie cosechada.
\[
Y_{\text{Curumaní, arroz, 2016}}
=\frac{131+25}{88+20}
=\frac{156}{108}
=1{,}44\text{ t/ha}.
\]
Al sumar las áreas antes de dividir, las 88 ha del primer semestre pesan más que las 20 ha
del segundo. La misma regla permite agregar municipios para construir las series departamentales.

\section{Atípicos con el rango intercuartílico}

El primer cuartil \(Q_1\) deja cerca del 25\,\% de los valores por debajo y el tercero \(Q_3\)
deja cerca del 75\,\%. La resta \(Q_3-Q_1\) mide la amplitud de la mitad central de los datos.
Para detectar extremos, usamos \(x=\ln(Y)\) en los rendimientos anuales positivos. En esta
escala, las distancias corresponden a razones entre rendimientos
\[
\ln(Y_2)-\ln(Y_1)=\ln(Y_2/Y_1).
\]
Subir de 1 a 2 t/ha y subir de 5 a 10 t/ha representa en ambos casos duplicar el
rendimiento. Las dos diferencias logarítmicas valen \(\ln(2)\). Esta propiedad motiva
el uso de distancias relativas dentro de cada cultivo. Calculamos las cercas del IQR
cuando el cultivo tiene al menos 20 rendimientos positivos y su IQR es mayor que cero.
\[
\mathrm{IQR}=Q_3-Q_1,\qquad
L=Q_1-3\,\mathrm{IQR},\qquad U=Q_3+3\,\mathrm{IQR}.
\]
Marcamos un extremo si \(x<L\) o \(x>U\). El factor 3 corresponde a las
\href{https://itl.nist.gov/div898/handbook/prc/section1/prc16.htm}{cercas exteriores}.
Los diagramas de caja suelen usar el factor 1,5 para sus bigotes.

El rendimiento anual de 1,44 t/ha que obtuvimos para Curumaní se compara con los demás registros de arroz.
En ese cultivo, los cuartiles del log-rendimiento son aproximadamente 1,534 y 1,736.
Con los valores sin redondear, el IQR es 0,203 y la cerca inferior es 0,926.
\[
\ln(156/108)\approx0{,}368 < 0{,}926
\qquad\Longleftrightarrow\qquad
1{,}44 < e^{0{,}926}\approx 2{,}52\text{ t/ha}.
\]
Curumaní en 2016 queda por debajo de esa cerca. De las 2.134 filas anuales, 11 tienen
rendimiento indefinido por falta de cosecha. El IQR se aplica a las 2.123 restantes
y marca \textbf{seis observaciones municipio, cultivo y año}. Se excluyen una de café,
una de algodón, una de arroz y tres de yuca.
\[
\begin{aligned}
2\,134-11 &= 2\,123 && \text{rendimientos definidos},\\
2\,123-6 &= 2\,117 && \text{rendimientos en el análisis depurado}.
\end{aligned}
\]

\begin{center}
\begin{minipage}{\linewidth}
\centering\small
\begin{tabular}{llrr}
\toprule
Cultivo & Municipio & Año & Rendimiento (t/ha)\\
\midrule
Café & Agustín Codazzi & 2011 & 0,20\\
Algodón & Chiriguaná & 2007 & 0,38\\
Arroz & Curumaní & 2016 & 1,44\\
Yuca & Río de Oro & 2016--2018 & 3,00 cada año\\
\bottomrule
\end{tabular}
\end{minipage}
\end{center}

Los seis casos anuales reúnen siete registros del CSV original porque Curumaní en 2016
aporta dos semestres. La tabla anual depurada tiene 2.128 filas. De ellas, las 11 sin
rendimiento se conservan para documentar la cobertura y quedan fuera de los cálculos.
La tabla completa permite repetir el análisis con los seis extremos. En Curumaní el arroz dio 4,47 t/ha
en 2015, 1,44 en 2016 y 4,35 en 2017. La yuca de Río de Oro pasó de 7 t/ha en 2015 a
3 t/ha durante tres años consecutivos con distintas superficies cosechadas.
Este segundo patrón es sostenido en los registros. Tratamos las exclusiones como una
decisión de sensibilidad y conservamos ambas versiones para interpretar los resultados.

\section{Forma y dispersión del rendimiento}
El diagrama de cajas muestra que cacao, café y fríjol suelen quedar por debajo de 1 t/ha,
mientras la yuca se concentra alrededor de 9 t/ha. Los productos se reportan en estados
físicos distintos, como grano seco, aceite crudo y tubérculo fresco. Las diferencias de nivel
entre cultivos reflejan también esa diferencia en lo que se pesa. Las comparaciones de
cambios se realizan dentro de cada cultivo.

\figura{.80\linewidth}{eva_R_02_cajas.png}{Rendimiento anual por cultivo en escala logarítmica.
Cada fila del gráfico reúne los rendimientos anuales de un cultivo. La línea interior es la mediana.
La caja va de \(Q_1\) a \(Q_3\). Los bigotes llegan a las observaciones más lejanas dentro
de las cercas de \(1{,}5\,\mathrm{IQR}\) y los puntos ocres destacan los extremos de \(3\,\mathrm{IQR}\).
El gráfico usa la tabla anual completa.}

\begin{center}
\begin{minipage}{\linewidth}
\centering\small
\begin{tabular}{lrrrrr}
\toprule
Cultivo & \(n\) & \(Q_1\) & Mediana & \(Q_3\) & DE\\
\midrule
Algodón & 92 & 1,64 & 2,20 & 2,50 & 0,67\\
Arroz & 215 & 4,64 & 5,03 & 5,69 & 0,79\\
Cacao & 219 & 0,40 & 0,50 & 0,60 & 0,11\\
Café & 227 & 0,56 & 0,62 & 0,76 & 0,16\\
Caña panelera & 89 & 3,00 & 4,00 & 4,50 & 1,05\\
Fríjol & 224 & 0,60 & 0,70 & 0,78 & 0,14\\
Maíz & 289 & 1,00 & 1,24 & 1,60 & 0,60\\
Palma de aceite & 241 & 3,00 & 3,30 & 3,58 & 0,51\\
Plátano & 236 & 5,00 & 6,10 & 7,00 & 1,65\\
Yuca & 285 & 8,00 & 9,03 & 10,00 & 1,81\\
\bottomrule
\end{tabular}
\captionof{table}{Descriptivos de la base depurada. Por cultivo, \(n\) cuenta combinaciones de
municipio y año con rendimiento definido, tras retirar los extremos. Cuartiles y desviación
estándar (DE) en t/ha.}
\end{minipage}
\end{center}

Los histogramas del análisis muestran una cola derecha visible en maíz y plátano,
y concentraciones repetidas en cacao y caña panelera. La mediana y los cuartiles permiten
resumir estas formas. Los histogramas por cultivo se adjuntan con las demás figuras.


\section{Variabilidad entre cultivos y dentro de cada cultivo}

En la base depurada, las diferencias entre las medias de los cultivos reúnen el
90,59\,\% de la variabilidad del rendimiento en t/ha. El cálculo asigna el mismo
peso a cada una de las 2.117 observaciones anuales con rendimiento definido.
Sea \(y_i\) el rendimiento de la observación \(i\), \(\bar y\) la media general,
\(\bar y_c\) la media del cultivo \(c\) y \(n_c\) su número de observaciones. La suma
de cuadrados total se descompone así
\[
\underbrace{\sum_i (y_i-\bar y)^2}_{\mathrm{Total}} =
\underbrace{\sum_c n_c(\bar y_c-\bar y)^2}_{\mathrm{Entre\ cultivos}} +
\underbrace{\sum_c\sum_{i\in c}(y_i-\bar y_c)^2}_{\mathrm{Dentro\ de\ los\ cultivos}}.
\]
La media general es 3,5317 t/ha. Por ejemplo, el cacao tiene 219 observaciones y una media
de 0,5119 t/ha. Su término en la suma entre cultivos es
\(219(0{,}5119-3{,}5317)^2\approx1\,997{,}1\).
Sumamos un término por cultivo y dividimos por la suma total.
\[
100\times\frac{19\,437{,}84}{21\,455{,}95}=90{,}59\,\%.
\]
El 9,41\,\% restante corresponde a diferencias
dentro de los cultivos. Con los seis extremos incluidos, la proporción entre cultivos
es 89,97\,\%. En escala logarítmica es 94,74\,\% en la base depurada.

\figura{.94\linewidth}{eva_Py_10_variabilidad.png}{Porcentaje de suma de cuadrados entre los grupos de cada variable.
El panel izquierdo usa la variabilidad total en t/ha como denominador. El derecho usa
la variabilidad que queda dentro de los cultivos. Las dos escalas horizontales son distintas.}

El municipio y el año definen sus propios grupos. Para el área sembrada, usamos los cortes de los quintiles.
En el panel derecho, cada observación se compara con la media de su propio cultivo.
Si \(g\) es el municipio, año o grupo de área dentro de \(c\), \(n_{cg}\) es su
número de registros y \(\bar y_{cg}\) es su media, la barra derecha se calcula como
\[
100\,\frac{\sum_c\sum_g n_{cg}(\bar y_{cg}-\bar y_c)^2}
{\sum_c\sum_{i\in c}(y_i-\bar y_c)^2}.
\]
Los cortes de área se calculan por separado en cada cultivo.

El municipio reúne la mayor fracción de la variabilidad interna en t/ha
(36,42\,\%), seguido del año (10,56\,\%) y los quintiles de área sembrada (6,93\,\%).
Estos porcentajes se calculan por separado y pueden solaparse. También dependen del número
de grupos, de su tamaño y, en el caso del área, de los cortes elegidos.
La identidad algebraica prescinde de supuestos de normalidad, independencia y varianza
constante. Las proporciones describen las diferencias de esta base.

\section{Cobertura y trayectoria anual}
La cobertura de datos condiciona la comparación. La yuca registra 24 municipios con
rendimiento definido en cada año y el café 19. El algodón llega a 14 municipios en 2011,
baja a dos en 2016 y carece de datos en 2018. La caña panelera tiene entre seis y ocho
municipios según el año. Esas diferencias anticipan qué cultivos tendrán evidencia limitada.

\figura{.90\linewidth}{eva_Py_06_cobertura.png}{Número de municipios con rendimiento definido
en la tabla completa de cada cultivo y año. Cero indica ausencia de observaciones disponibles.}

\figura{.93\linewidth}{eva_R_04_trayectorias.png}{Rendimiento departamental ponderado por área cosechada.
Las escalas verticales son independientes. El arroz cae durante buena parte del periodo y repunta
en 2018. La palma baja después de 2012 y recupera parte del nivel al final.}

Cada punto de la trayectoria agrega los municipios de un cultivo en un año. Por ejemplo, el
punto de arroz en 2018 combina 46.247 t en 7.722 ha cosechadas y vale
\(46\,247/7\,722=5{,}99\) t/ha. La cosecha total pondera el aporte de cada municipio.

\section{Diferencias municipales}
Para comparar territorios, restamos a cada rendimiento la mediana de su cultivo y año.
En maíz, San Martín quedó 1,82 t/ha por encima de esa referencia en la mediana de las
desviaciones de 12 años. El Paso quedó 0,31 t/ha por debajo con la misma cobertura.
Para construir el mapa, expresamos la desviación de cada año en porcentaje de la mediana de
referencia \(M_{ct}\). Cada celda resume la mediana de los porcentajes de un municipio y cultivo.
\[
D_{mct}=100\,\frac{Y_{mct}-M_{ct}}{M_{ct}}.
\]
Un valor de 20\,\% indica un rendimiento un quinto por encima de la referencia.

\figura{.64\linewidth}{eva_Py_07_territorio.png}{Mediana municipal de las desviaciones porcentuales
respecto a la mediana del mismo cultivo y año. La escala de color satura en \(\pm40\,\%\).
La tabla territorial adjunta indica cuántos años sustentan cada celda.}

\section{Conclusiones}

La mayor parte de la variabilidad observada del rendimiento corresponde a las diferencias
entre las medias de los cultivos. En t/ha, esa fracción es 90,59\,\% con la base depurada y
89,97\,\% al incluir los seis extremos. En la comparación interviene el estado físico en que
EVA registra cada producto. Una tonelada de grano seco y una tonelada de tubérculo fresco
representan productos distintos.

Al comparar registros del mismo cultivo, la separación por municipio reúne el 36,42\,\%
de la variabilidad interna en t/ha. La separación por año reúne el 10,56\,\% y por grupos
de área sembrada, el 6,93\,\%. Cada cálculo compara grupos de una sola variable, de modo que
las fracciones pueden solaparse. El maíz ilustra la diferencia territorial con San Martín
por encima de la referencia de su cultivo y El Paso por debajo, ambos con 12 años observados.

La cobertura condiciona la lectura de las series departamentales. La yuca tiene 24 municipios
con rendimiento definido en cada año. El algodón pasa de 14 en 2011 a dos en 2016 y a ninguno
en 2018. Estas series resumen la producción y la cosecha de los municipios disponibles
en cada año. Los cambios observados pueden combinar diferencias de rendimiento y de cobertura.

\subsection*{Reproducibilidad}
La carpeta del proyecto contiene el CSV de EVA, scripts de R, cuadernos de Colab, tablas y
figuras. La ejecución local contrastó filas, atípicos y descriptivos entre ambos lenguajes. Las cifras del informe se cotejaron con las salidas de ambos flujos.
La revisión visual comprobó los ejes, las unidades y la correspondencia entre cada figura
y su interpretación.

\paragraph{Uso de IA}
Se usó IA generativa como apoyo para escribir el código.

\end{document}
